% Aqui começa o capítulo de Planejamento Ágil.
\chapter{Planejamento Ágil}\label{chp:PlanejamentoAgil}

O planejamento ágil do projeto de desenvolvimento do sistema de gestão de casos sociais é orientado pelos princípios de entrega incremental de valor, adaptação contínua e transparência. Em consonância com o PMBOK 7ª edição, o planejamento não é tratado como uma atividade estática realizada apenas no início do projeto, mas como um processo contínuo que evolui à medida que novas informações são obtidas.

Neste capítulo são apresentados os principais elementos do planejamento ágil adotado, incluindo a visão do produto, a estrutura do backlog, o roadmap de entregas e as métricas utilizadas para acompanhamento e controle do projeto.

\section{Visão do Produto}

A visão do produto estabelece um entendimento compartilhado sobre o propósito e o valor do sistema a ser desenvolvido. O sistema de gestão de casos sociais tem como objetivo central apoiar equipes técnicas e gestores na organização, acompanhamento e análise de casos sociais, promovendo maior rastreabilidade das informações, integração territorial e apoio à tomada de decisão.

A solução é concebida como uma aplicação digital que permita o cadastro estruturado de casos, a categorização por tags temáticas, o registro evolutivo de atendimentos, a visualização georreferenciada em mapa e a geração de resumos analíticos. Essa visão orienta a priorização do trabalho e serve como referência para decisões ao longo de todo o projeto.

\section{Estrutura do Product Backlog}

O Product Backlog representa a lista priorizada de funcionalidades, requisitos e melhorias do sistema, organizada de acordo com o valor entregue aos stakeholders. No contexto deste projeto, o backlog é estruturado em níveis, iniciando por épicos que representam grandes blocos funcionais, os quais são progressivamente refinados em histórias de usuário mais granulares.

Os principais épicos definidos para o sistema são:
\begin{itemize}
    \item Gestão de Casos Sociais.
    \item Categorização e Organização por Tags.
    \item Registro de Apontamentos e Evolução dos Casos.
    \item Visualização Georreferenciada em Mapa.
    \item Relatórios e Resumos Analíticos.
    \item Gestão de Usuários e Perfis de Acesso.
\end{itemize}

Cada épico é desdobrado em histórias de usuário que descrevem funcionalidades sob a perspectiva dos usuários finais, acompanhadas de critérios de aceitação claros. O backlog é continuamente refinado ao longo do projeto, incorporando feedback dos stakeholders e aprendizados obtidos a cada sprint.

\section{Planejamento das Sprints}

O planejamento das sprints ocorre no início de cada ciclo, no qual a equipe define os objetivos da sprint e seleciona os itens do Product Backlog que serão trabalhados, considerando sua capacidade e as prioridades estabelecidas pelo Product Owner.

As sprints possuem duração fixa, o que contribui para a previsibilidade e facilita a comparação de desempenho ao longo do tempo. Ao final de cada sprint, espera-se a entrega de um incremento funcional do sistema, potencialmente utilizável ou avaliável pelos stakeholders.

Esse modelo de planejamento permite manter controle sobre o progresso do projeto, ao mesmo tempo em que preserva flexibilidade para ajustes de escopo e prioridade.

\section{Roadmap do Projeto}

O roadmap fornece uma visão de médio prazo das entregas previstas, sem detalhamento excessivo, reconhecendo a existência de incertezas. No projeto em análise, o roadmap é organizado por conjuntos de sprints, associados a incrementos de valor claramente identificáveis.

De forma geral, o roadmap contempla:
\begin{itemize}
    \item Fase inicial focada na infraestrutura básica do sistema e no cadastro de casos.
    \item Fase intermediária dedicada à organização por tags, registro de apontamentos e melhorias de usabilidade.
    \item Fase avançada voltada à visualização em mapa, geração de resumos analíticos e consolidação do produto.
\end{itemize}

O roadmap é revisado periodicamente, podendo ser ajustado conforme mudanças de contexto, feedback dos usuários ou restrições institucionais. Essa abordagem permite alinhar expectativas dos stakeholders sem comprometer a adaptabilidade do projeto.

\section{Métricas Ágeis e Acompanhamento}

A medição do desempenho do projeto é realizada por meio de métricas ágeis simples, alinhadas aos objetivos do projeto e aos domínios de desempenho do PMBOK 7. O uso de métricas busca apoiar a tomada de decisão e o aprendizado contínuo, evitando controles excessivamente burocráticos.

Entre as métricas adotadas destacam-se:
\begin{itemize}
    \item Progresso do Product Backlog, por meio do acompanhamento dos itens concluídos.
    \item Previsibilidade das sprints, observando a relação entre trabalho planejado e trabalho entregue.
    \item Estabilidade do ritmo de trabalho da equipe ao longo do tempo.
    \item Feedback qualitativo dos stakeholders durante as revisões de sprint.
\end{itemize}

Essas métricas são analisadas de forma conjunta, considerando o contexto do projeto e evitando interpretações isoladas. O foco principal é identificar tendências, apoiar ajustes no planejamento e promover a melhoria contínua do processo de trabalho.

\section{Previsibilidade e Controle}

Embora o projeto adote uma abordagem ágil, a previsibilidade é tratada como um objetivo relevante, especialmente em um contexto institucional. A combinação de sprints de duração fixa, backlog priorizado e métricas de acompanhamento permite manter controle sobre o andamento do projeto e fornecer informações claras aos stakeholders.

Dessa forma, o planejamento ágil adotado busca equilibrar flexibilidade e controle, reconhecendo as incertezas inerentes ao projeto, mas mantendo alinhamento com os objetivos estratégicos e os princípios do PMBOK 7ª edição.
